\documentclass[12pt,a4paper]{article}

% Packages
\usepackage{graphicx}
\usepackage{geometry}
\usepackage{setspace}
\usepackage{float}

\geometry{margin=1in}
\onehalfspacing

%-----------------------------------------
% Project Information
%-----------------------------------------

\title{
    \textbf{Digital Clock Project}\\[0.5em]
    \large Progress Log
}

\author{
    \textbf{Member 1:} G.N.S.S.Rithvik\\
    Roll Number: ee26btech11211
    \\[0.5em]
    \textbf{Member 2:} R.Srikar Varshit\\
    Roll Number: ee26btech11223
}

\date{\today}

%-----------------------------------------

\begin{document}

\maketitle

\tableofcontents
\newpage


%=========================================
\section{Project Start}
%=========================================

We started the Digital Clock project by understanding the requirements and discussing how the system could be implemented. Both members contributed to the initial planning and discussion.

\subsection*{G.Rithvik}

\begin{itemize}
    \item Worked on the programming part that is required for the clock.
    \item Helped in construction the circuit.
    \item Contributed to the initial circuit planning.
\end{itemize}

\subsection*{R.Srikar}

\begin{itemize}
    \item Helped understand the project requirements.
    \item Researched about the apparatus and connections of circuit.
    \item Contributed to documenting the whole project.
\end{itemize}

\begin{figure}[h]
    \centering
    \includegraphics[width=0.7\textwidth, angle=90]{digital clock/digi c origin.jpeg}
    \caption{Initial project setup and planning}
\end{figure}


%=========================================
\section{Initial Circuit Setup}
%=========================================

We began assembling the Digital Clock circuit according to our
planned design. We checked the connections and tested the circuit
step by step.

\subsection*{G.Rithvik}

\begin{itemize}
    \item Worked on assembling the circuit.
    \item Checked the connections between the components.
    \item Helped test the initial circuit.
\end{itemize}

\subsection*{R.Srikar}

\begin{itemize}
    \item Assisted with circuit assembly.
    \item Checked the power and component connections.
    \item Helped identify and correct connection issues.
\end{itemize}

\textbf{Problem faced:} We got a problem with the displays.They did not show the digits properly

\textbf{Solution:} It was the connection issue.we reconnected the jumpers again.

\begin{figure}[h]
    \centering
    \includegraphics[width=0.7\textwidth]{digital clock/digi c constructed.jpeg}
    \caption{Initil circuit assembly}
\end{figure}


%=========================================
\section{Programming and Testing}
%=========================================

After assembling the circuit, we worked on the Arduino program
and tested it with the hardware.

\subsection*{G.Rithvik}

\begin{itemize}
    \item Worked on the Arduino code.
    \item Tested the program with the circuit.
    \item Helped debug errors in the program.
\end{itemize}

\subsection*{R.Srikar}

\begin{itemize}
    \item Assisted in testing the program.
    \item Checked whether the displayed time was correct.
    \item Helped identify problems during testing.
\end{itemize}

\textbf{What we learned:}

K-Maps, State Machines, Multiplexing, Circuit Building, Code

\begin{figure}[h]
    \centering
    \includegraphics[width=0.7\textwidth,angle=90]{digital clock/digi c error.jpeg}
    \caption{Testing the Digital Clock}
\end{figure}


%=========================================   
\section{Debugging}
%=========================================

During testing, we encountered some problems with the circuit. We worked together to identify the causes and
make the required changes.

\subsection*{Problem}

The problem we encountered was the displays were not displaying the digits properly.There were also some sections that wouldn't get power/light up.

\subsection*{Our Approach}

Our approach to our problem we encountered is we unpluging all the jumper wires and checking whether all the displays are working with a working jumper  wire.There were some broken jumper wires we removed them.

\subsection*{Solution}

We connected the jumper wires in a proper manner to the breadboard.we figured out that the issue was loose connection of jumper wires and some faulty jumper wires


%=========================================
\section{Final Testing}
%=========================================

After making the required changes, we tested the complete
Digital Clock. Both members checked the operation of the system
and verified the output.

\subsection*{G.Rithvik}

verified the arduino code.

\subsection*{R.Srikar}

Checked whether the displayed time was correct.

\textbf{Labelling of the push buttons:}
\begin{itemize}
    \item from left to right the push buttons are labelled as {0,1,2,3}
\end{itemize}
\textbf{Function of push buttons:}
\begin{itemize}
    \item $0\rightarrow$ Start/Pause
    \item $1\rightarrow$ Next
    \item $2\rightarrow$ Increment
    \item $3\rightarrow$ Decrement
\end{itemize}


\begin{figure}[h]
    \centering
    \includegraphics[width=0.7\textwidth,angle=90]{digital clock/digi clock working.jpeg}
    \caption{Final working Digital Clock}
\end{figure}


%=========================================
\section{Individual Learning}
%=========================================

\subsection*{G.Rithvik}

Through this project, I learned about the implementation and
working of a Digital Clock. I also learned about
\textbf {K-Maps and its uses, circuit assembling, uses of components used in the making of digital clock.}

The most difficult part for me was
\textbf {Writing and understanding the code which is used in the hardware}

The most interesting part was
\textbf{Figuring out the connections made between the components in the circuit}


\subsection*{R.Srikar}

Through this project, I learned about the implementation and
working of a Digital Clock. I also learned about
\textbf {K-Maps and its uses, circuit assembling, uses of components used in the making of digital clock.}

The most difficult part for me was
\textbf {Figuring out the mistakes made in connecting the components}

The most interesting part was
\textbf{Figuring out the connections made between the components in the circuit}
 
\begin{figure}[h]
    \centering
    \includegraphics[width=1.1\textwidth,]{digital clock/K-Maps.png}
    \caption{K-Maps of the code}
\end{figure}

 
%=========================================
\section{Final Reflection}
%=========================================

Working as a two-member team allowed us to divide the work and
solve problems together. The project involved planning,
assembling the circuit, programming, testing, debugging, and
final verification.



\end{document}
```
